% GENERATED FILE. Edit canonical lessons, then run npm run generate:books.
% canonical-source-hash: fnv1a64:e4a25a05ab6af0f3
% canonical-lessons: ML-C440-vyavastha, ML-C440-nibandhana, ML-C440-upadhi, ML-C440-illattapaksam, ML-C440-allattapaksam

\chapter{On One Condition}
\label{ch:ml-on-one-condition}
\hlchaptermodality{440}

\begin{chapteropening}
\textbf{By the end of this chapter:} \emph{I can say a condition, a stipulation, a proviso, if not and otherwise.}

Complete the last lesson of chapter 440: I can say a condition, a stipulation, a proviso, if not and otherwise.
\end{chapteropening}

\section[\texorpdfstring{\ml{വ്യവസ്ഥ} (vyavastha)}{വ്യവസ്ഥ (vyavastha)}]{\ml{വ്യവസ്ഥ} (vyavastha) — a condition; a system}

\label{lesson:ML-C440-vyavastha}



\begin{quote}
Before the new one: say the Malayalam for a reference, then the Malayalam for an author.
\end{quote}

\subsection*{You'll want to know: \ml{വ്യവസ്ഥ}}
\textbf{\ml{വ്യവസ്ഥ}} — \emph{vyavastha} — \textquotedblleft{}a condition; a system\textquotedblright{}.

\begin{quote}
\ml{നീയും} \ml{വന്നാൽ} \ml{മാത്രം} \ml{ഞാൻ} \ml{വരാം}, \ml{അതാണ്} \ml{എന്റെ} \ml{വ്യവസ്ഥ}.
\end{quote}

\begin{quote}
\emph{nīyuṁ vannāl mātraṁ ñān varāṁ, atāṇŭ enṟe vyavastha.}
\end{quote}

\begin{quote}
I will come only if you come too; that is my condition.
\end{quote}

It is also an order of things, a system: \ml{സാമ്പത്തിക} \ml{വ്യവസ്ഥ}, the economic system.

\begin{cousinweb}
\textbf{\ml{വ്യവസ്ഥ}} (\emph{vyavastha}) is borrowed from Sanskrit \emph{vyavasthā}, from the verb \emph{vy-ava-sthā}, \textquotedblleft{}to stand apart; to be settled, fixed\textquotedblright{} (Monier-Williams): a condition is a thing settled in advance.
\end{cousinweb}

\begin{grammarlens}[title={What you've built}]
One more word for saying what happens if.
\end{grammarlens}

\subsection*{Your turn}
\begin{itemize}
\raggedright
  \item \emph{Say it:} \emph{vyavastha}
  \item \emph{Say it:} \emph{vyavastha}, once more
  \item \emph{Say it:} \emph{nīyuṁ vannāl mātraṁ ñān varāṁ, atāṇŭ enṟe vyavastha.} — I will come only if you come too; that is my condition.
  \item \emph{Recall:} say the Malayalam for a reference, then the Malayalam for an author, then say \emph{vyavastha} again
\end{itemize}

\begin{tcolorbox}[breakable,colback=teal!4,colframe=teal!35!black,title={Before you move on}]
What does \ml{വ്യവസ്ഥ} mean? (\textquotedblleft{}A condition; a system\textquotedblright{}.) Say it once more.
\end{tcolorbox}
\section[\texorpdfstring{\ml{നിബന്ധന} (nibandhana)}{നിബന്ധന (nibandhana)}]{\ml{നിബന്ധന} (nibandhana) — a stipulation; a term}

\label{lesson:ML-C440-nibandhana}



\begin{quote}
Before the new one: say the Malayalam for an author, then the Malayalam for a condition.
\end{quote}

\subsection*{You'll want to know: \ml{നിബന്ധന}}
\textbf{\ml{നിബന്ധന}} — \emph{nibandhana} — \textquotedblleft{}a stipulation; a term\textquotedblright{}.

\begin{quote}
\ml{കരാറിൽ} \ml{പല} \ml{നിബന്ധനകളും} \ml{ഉണ്ട്}.
\end{quote}

\begin{quote}
\emph{karāṟil pala nibandhanakaḷuṁ uṇṭŭ.}
\end{quote}

\begin{quote}
The contract has many stipulations.
\end{quote}

Plural: \ml{നിബന്ധനകൾ}, \emph{nibandhanakaḷ}.

Plural: \ml{നിബന്ധനകൾ}, the terms of an agreement.

\begin{cousinweb}
\textbf{\ml{നിബന്ധന}} (\emph{nibandhana}) is borrowed from Sanskrit \emph{nibandhana}, \textquotedblleft{}binding, fastening\textquotedblright{} (Monier-Williams): a stipulation binds the two sides.
\end{cousinweb}

\begin{grammarlens}[title={What you've built}]
One more word for saying what happens if.
\end{grammarlens}

\subsection*{Your turn}
\begin{itemize}
\raggedright
  \item \emph{Say it:} \emph{nibandhana}
  \item \emph{Say it:} \emph{nibandhana}, once more
  \item \emph{Say it:} \emph{karāṟil pala nibandhanakaḷuṁ uṇṭŭ.} — The contract has many stipulations.
  \item \emph{Recall:} say the Malayalam for an author, then the Malayalam for a condition, then say \emph{nibandhana} again
\end{itemize}

\begin{tcolorbox}[breakable,colback=teal!4,colframe=teal!35!black,title={Before you move on}]
What does \ml{നിബന്ധന} mean? (\textquotedblleft{}A stipulation; a term\textquotedblright{}.) Say it once more.
\end{tcolorbox}
\section[\texorpdfstring{\ml{ഉപാധി} (upādhi)}{ഉപാധി (upādhi)}]{\ml{ഉപാധി} (upādhi) — a proviso; a condition}

\label{lesson:ML-C440-upadhi}



\begin{quote}
Before the new one: say the Malayalam for a condition, then the Malayalam for a stipulation.
\end{quote}

\subsection*{You'll want to know: \ml{ഉപാധി}}
\textbf{\ml{ഉപാധി}} — \emph{upādhi} — \textquotedblleft{}a proviso; a condition\textquotedblright{}.

\begin{quote}
\ml{ഉപാധികൾ} \ml{ഒന്നുമില്ലാതെ} \ml{അവർ} \ml{സഹായം} \ml{നൽകി}.
\end{quote}

\begin{quote}
\emph{upādhikaḷ onnumillāte avar sahāyaṁ nalki.}
\end{quote}

\begin{quote}
They gave help without any conditions.
\end{quote}

Unconditional is \ml{ഉപാധികളില്ലാത്ത}.

\begin{cousinweb}
\textbf{\ml{ഉപാധി}} (\emph{upādhi}) is borrowed from Sanskrit \emph{upādhi}, \textquotedblleft{}a point of view, an aim\textquotedblright{} (CDIAL 2299), and in Monier-Williams \textquotedblleft{}that which is put in place of another thing\textquotedblright{}.
\end{cousinweb}

\begin{grammarlens}[title={What you've built}]
One more word for saying what happens if.
\end{grammarlens}

\subsection*{Your turn}
\begin{itemize}
\raggedright
  \item \emph{Say it:} \emph{upādhi}
  \item \emph{Say it:} \emph{upādhi}, once more
  \item \emph{Say it:} \emph{upādhikaḷ onnumillāte avar sahāyaṁ nalki.} — They gave help without any conditions.
  \item \emph{Recall:} say the Malayalam for a condition, then the Malayalam for a stipulation, then say \emph{upādhi} again
\end{itemize}

\begin{tcolorbox}[breakable,colback=teal!4,colframe=teal!35!black,title={Before you move on}]
What does \ml{ഉപാധി} mean? (\textquotedblleft{}A proviso; a condition\textquotedblright{}.) Say it once more.
\end{tcolorbox}
\section[\texorpdfstring{\ml{ഇല്ലാത്തപക്ഷം} (illāttapakṣaṁ)}{ഇല്ലാത്തപക്ഷം (illāttapakṣaṁ)}]{\ml{ഇല്ലാത്തപക്ഷം} (illāttapakṣaṁ) — if not; failing that}

\label{lesson:ML-C440-illattapaksam}



\begin{quote}
Before the new one: say the Malayalam for a stipulation, then the Malayalam for a proviso.
\end{quote}

\subsection*{You'll want to know: \ml{ഇല്ലാത്തപക്ഷം}}
\textbf{\ml{ഇല്ലാത്തപക്ഷം}} — \emph{illāttapakṣaṁ} — \textquotedblleft{}if not; failing that\textquotedblright{}.

\begin{quote}
\ml{ഫീസ്} \ml{ഇന്ന്} \ml{അടയ്ക്കണം}, \ml{ഇല്ലാത്തപക്ഷം} \ml{പേര്} \ml{വെട്ടും}.
\end{quote}

\begin{quote}
\emph{phīsŭ innŭ aṭaykkaṇaṁ, illāttapakṣaṁ pērŭ veṭṭuṁ.}
\end{quote}

\begin{quote}
The fee must be paid today; if not, the name will be struck off.
\end{quote}

A formal partner of \ml{ഇല്ലെങ്കിൽ}.

\begin{cousinweb}
\textbf{\ml{ഇല്ലാത്തപക്ഷം}} (\emph{illāttapakṣaṁ}) is a Malayalam compound: \ml{ഇല്ലാത്ത}, \textquotedblleft{}that is not\textquotedblright{}, native Dravidian (DEDR 2559), and \ml{പക്ഷം}, \textquotedblleft{}a side; a case\textquotedblright{}, from Sanskrit \emph{pakṣa} (CDIAL 7627): in the case that it is not.
\end{cousinweb}

\begin{grammarlens}[title={What you've built}]
One more word for saying what happens if.
\end{grammarlens}

\subsection*{Your turn}
\begin{itemize}
\raggedright
  \item \emph{Say it:} \emph{illāttapakṣaṁ}
  \item \emph{Say it:} \emph{illāttapakṣaṁ}, once more
  \item \emph{Say it:} \emph{phīsŭ innŭ aṭaykkaṇaṁ, illāttapakṣaṁ pērŭ veṭṭuṁ.} — The fee must be paid today; if not, the name will be struck off.
  \item \emph{Recall:} say the Malayalam for a stipulation, then the Malayalam for a proviso, then say \emph{illāttapakṣaṁ} again
\end{itemize}

\begin{tcolorbox}[breakable,colback=teal!4,colframe=teal!35!black,title={Before you move on}]
What does \ml{ഇല്ലാത്തപക്ഷം} mean? (\textquotedblleft{}If not; failing that\textquotedblright{}.) Say it once more.
\end{tcolorbox}
\section[\texorpdfstring{\ml{അല്ലാത്തപക്ഷം} (allāttapakṣaṁ)}{അല്ലാത്തപക്ഷം (allāttapakṣaṁ)}]{\ml{അല്ലാത്തപക്ഷം} (allāttapakṣaṁ) — or else; otherwise}

\label{lesson:ML-C440-allattapaksam}



\begin{quote}
Before the new one: say the Malayalam for a proviso, then the Malayalam for if not.
\end{quote}

\subsection*{You'll want to know: \ml{അല്ലാത്തപക്ഷം}}
\textbf{\ml{അല്ലാത്തപക്ഷം}} — \emph{allāttapakṣaṁ} — \textquotedblleft{}or else; otherwise\textquotedblright{}.

\begin{quote}
\ml{അവൻ} \ml{വേഗം} \ml{ഇറങ്ങി}, \ml{അല്ലാത്തപക്ഷം} \ml{ബസ്} \ml{പോകുമായിരുന്നു}.
\end{quote}

\begin{quote}
\emph{avan vēgaṁ iṟaṅṅi, allāttapakṣaṁ basŭ pōkumāyirunnu.}
\end{quote}

\begin{quote}
He left quickly; otherwise the bus would have gone.
\end{quote}

It joins a clause and the bad result that would follow without it.

\begin{cousinweb}
\textbf{\ml{അല്ലാത്തപക്ഷം}} (\emph{allāttapakṣaṁ}) is a Malayalam compound: \ml{അല്ലാത്ത}, \textquotedblleft{}that is not so\textquotedblright{}, native Dravidian (DEDR 234), and \ml{പക്ഷം}, \textquotedblleft{}a case\textquotedblright{}, from Sanskrit \emph{pakṣa} (CDIAL 7627).
\end{cousinweb}

\begin{grammarlens}[title={What you've built}]
One more word for saying what happens if.
\end{grammarlens}

\subsection*{Your turn}
\begin{itemize}
\raggedright
  \item \emph{Say it:} \emph{allāttapakṣaṁ}
  \item \emph{Say it:} \emph{allāttapakṣaṁ}, once more
  \item \emph{Say it:} \emph{avan vēgaṁ iṟaṅṅi, allāttapakṣaṁ basŭ pōkumāyirunnu.} — He left quickly; otherwise the bus would have gone.
  \item \emph{Recall:} say the Malayalam for a proviso, then the Malayalam for if not, then say \emph{allāttapakṣaṁ} again
\end{itemize}

\begin{tcolorbox}[breakable,colback=teal!4,colframe=teal!35!black,title={Before you move on}]
What does \ml{അല്ലാത്തപക്ഷം} mean? (\textquotedblleft{}Or else; otherwise\textquotedblright{}.) Say it once more.
\end{tcolorbox}
